\documentclass[10pt,a4paper]{amsart}
\usepackage[margin=19mm]{geometry}
\usepackage{amsmath,amssymb,mathtools}
\usepackage[scr=rsfs,cal=euler]{mathalfa}
\usepackage{xfrac}
\usepackage[unicode,hidelinks,bookmarks=false]{hyperref}
\newcommand{\ie}{i.e.\ }
\newcommand{\cf}{cf.\ }
\newcommand{\resp}{respectively\ }
\newcommand{\Subsection}{\subsection}
\providecommand{\nobreakdash}{\nobreak\hbox{-}}
\setlength{\emergencystretch}{2em}
\multlinegap0pt
\usepackage{amssymb}
\usepackage{mathtools}
\usepackage[shortlabels]{enumitem}
\newtheorem{lemma}{Lemma}
\newcommand{\artanh}{\operatorname{artanh}}
\title{The Probability That the Incenter of a\\Triangle Lies in a Random Diameter Disk}
\author{Stanley Rabinowitz}
\date{Source: arXiv:2608.27491v1 (CC BY 4.0)}
\begin{document}
\maketitle

\noindent\emph{Source and licence.} The Probability That the Incenter of a\\Triangle Lies in a Random Diameter Disk. Version arXiv:2608.27491v1 (CC BY 4.0), distributed by arXiv under the Creative Commons Attribution 4.0 International licence, http://creativecommons.org/licenses/by/4.0/, as recorded in the arXiv licence metadata for this exact version and shown on the abstract page https://arxiv.org/abs/2608.27491v1. Full licence notice: \texttt{licenses/arxiv-2608.27491-CC-BY-4.0.txt}.\par\smallskip

\noindent\textit{Editorial scope.} Counted unit: the article's lemma lem:H4 (source0.tex lines 854--870). The article's complete original proof of it is delivered with it (source0.tex lines 872--988, one \texttt{proof} environment of 117 lines). No mathematics is written, repaired or re-derived: the statement and the proof are the article's own lines, in the article's own notation, and no retained line is substituted or re-flowed.  Retained uncounted as the source-local context the proof needs: lines 478--481; lines 818--827.  Nothing was omitted from any retained span, so the package carries no omission record.  No retained line contains a citation, so the wrapper carries no bibliography and no bibliography entry.  No retained line contains a figure environment, an image-inclusion command or drawing code, so no image is delivered and the package declares no asset dependency.  Editorial formatting only, in addition: an AMS \texttt{amsart} preamble that restates the article's own declarations and macros used by the retained lines, namely the environments \texttt{lemma} on separate counters, together with the standard packages those lines need.  The retained environments print in this document as Lemma 1 in order of appearance, whereas the article prints the same items with the numbering of its own full document.  The complete native source of the article as distributed in the arXiv e-print for this exact version is bundled unchanged as \texttt{sources/208717/source0.tex}.
\begin{equation}
\label{eq:kappa}
\kappa=10-9\log3.
\end{equation}

\begin{lemma}
\label{lem:atanhbound}
For $0<t<1$,
\begin{equation}
\label{eq:atanhbound}
\artanh t
>
\frac{t(945-735t^2+64t^4)}{15(63-70t^2+15t^4)}.
\end{equation}
\end{lemma}

\begin{lemma}
\label{lem:H4}
Let
\begin{equation}
\label{eq:Hdef}
H(u)
=
\Gamma(1-2u^2)+2\Gamma(u)-1-2\kappa u(1-u),
\qquad 0<u<1,
\end{equation}
where $\kappa$ is given by \eqref{eq:kappa}.  Then
\begin{equation}
\label{eq:H4neg}
H^{(4)}(u)<0
\qquad(0<u<1).
\end{equation}
\end{lemma}

\begin{proof}
The quadratic term involving $\kappa$ disappears after four differentiations.  Put
\[
a=1-2u^2.
\]
For $u\ne1/\sqrt2$, direct differentiation gives
\begin{align}
H^{(4)}(u)={}&
\frac{8N_{16}(u)}
{u^5(u-1)^3(u+1)^3(2u^2-1)^5}
-\frac{240}{u^6}\artanh u \notag\\
&-\frac{96(140u^4+84u^2+3)}{(2u^2-1)^6}\artanh a,
\label{eq:H4exact}
\end{align}
where
\begin{align*}
N_{16}(u)={}&384u^{16}-480u^{15}-3072u^{14}+960u^{13}+8800u^{12}
+564u^{11}\\
&-13120u^{10}-1884u^9+11560u^8+743u^7-6272u^6+106u^5\\
&+2066u^4-u^3-380u^2+30.
\end{align*}
The apparent singularity at $u=1/\sqrt2$ is removable, since $H$ is smooth there.  Formula~\eqref{eq:H4exact} makes clear, in particular, that the coefficients of $\artanh u$ and $\artanh a$ are both negative.

First suppose $0<u<1/\sqrt2$, so $0<a<1$.  Apply Lemma~\ref{lem:atanhbound} to both inverse-hyperbolic-tangent terms.  Because both coefficients are negative, this gives an upper bound.
Define
\[
D_1(u)=(15u^4-70u^2+63)
(30u^8-60u^6+10u^4+20u^2+1).
\]
Then the resulting upper bound is
\begin{equation}
\label{eq:H4case1}
H^{(4)}(u)
\le
\frac{8P_{15}(u)}{5u^2(u^2-1)^3D_1(u)},
\end{equation}
where
\begin{align*}
P_{15}(u)={}&
7800u^{15}-150u^{14}-12000u^{13}-12515u^{12}
-20800u^{11}+78035u^{10}\\
&+38800u^9-138383u^8-2740u^7+88504u^6
-10680u^5-17089u^4\\
&-540u^3+1603u^2+315.
\end{align*}
This polynomial is positive for $0<u<1$.  Indeed, with $u=y/(1+y)$,
\begin{align*}
&(1+y)^{15}P_{15}\left(\frac{y}{1+y}\right)\\
={}&315+4725y+34678y^2+163624y^3+531440y^4
+1170104y^5\\
&+1645729y^6+1316297y^7+454964y^8+118180y^9
+397704y^{10}\\
&+494324y^{11}+279824y^{12}+93360y^{13}+8320y^{14}+160y^{15},
\end{align*}
whose coefficients are all positive.  The two factors defining $D_1(u)$ are positive on the present interval; in fact
\[
15u^4-70u^2+63>0
\]
and
\[
30u^8-60u^6+10u^4+20u^2+1
=\frac18\bigl(15a^4-70a^2+63\bigr)>0.
\]
Since $(u-1)^3(u+1)^3<0$, the right side of \eqref{eq:H4case1} is negative.

Now suppose $1/\sqrt2<u<1$ and put
\[
b=2u^2-1=-a\in(0,1).
\]
Then $\artanh a=-\artanh b$.  We use the elementary bounds
\begin{equation}
\label{eq:atanhlowertrunc}
\artanh t>
 t+\frac{t^3}{3}+\frac{t^5}{5}+\frac{t^7}{7}
\end{equation}
and
\begin{equation}
\label{eq:atanhuppertrunc}
\artanh t<
 t+\frac{t^3}{3}+\frac{t^5}{5}
 +\frac{t^7}{7(1-t^2)},
\qquad 0<t<1.
\end{equation}
The first follows by truncating the positive series for $\artanh t$; the second follows by bounding its remaining denominators by $7$.

Apply \eqref{eq:atanhlowertrunc} to $\artanh u$ and \eqref{eq:atanhuppertrunc} to $\artanh b$.  The result simplifies to
\begin{equation}
\label{eq:H4case2}
H^{(4)}(u)
\le
\frac{16P_{10}(u)}{35u^2(u-1)^3(u+1)^3},
\end{equation}
where
\begin{align*}
P_{10}(u)={}&
840u^{10}-75u^9-1596u^8+120u^7+928u^6
-85u^5\\
&+102u^4-30u^3-174u^2+40.
\end{align*}
This polynomial is positive for $u>3/5$.  Put
\[
u=\frac{3/5+y}{1+y},
\qquad y>0.
\]
Then
\begin{align*}
&5^9(1+y)^{10}P_{10}\left(\frac{3/5+y}{1+y}\right)\\
={}&3246602+16762640y+93221670y^2+536913600y^3
+1813552500y^4\\
&+3534460000y^5+4220537500y^6+3261000000y^7
+1738281250y^8\\
&+656250000y^9+136718750y^{10},
\end{align*}
again with all coefficients positive.  Since $1/\sqrt2>3/5$ and $(u-1)^3(u+1)^3<0$, the right side of \eqref{eq:H4case2} is negative.

The point $u=1/\sqrt2$ follows by continuity (or by taking the limit in either bound).  This proves \eqref{eq:H4neg}.
\end{proof}

\end{document}
